\chapter{Introductory Concepts}
\label{chap:introductory_concepts}

This chapter introduces the central object of study: a \textbf{metric space}. We examine canonical examples, sphere geometries, open and closed sets, interior, closure, boundary, distance between sets, diameter, subspaces, and product spaces.

\section{Definition and Examples of Metric Spaces}
\label{sec:metric_space_def}

\begin{definition}[Metric Space]
Let $X$ be a non-empty set. A \textbf{metric} (or distance function) on $X$ is a real-valued function $d: X \times X \to \R$ satisfying the following axioms for all $x, y, z \in X$:
\begin{enumerate}[label=\rm{(M\arabic*)}]
    \item \textbf{Non-negativity}: $d(x,y) \ge 0$.
    \item \textbf{Identity of Indiscernibles}: $d(x,y) = 0 \iff x = y$.
    \item \textbf{Symmetry}: $d(x,y) = d(y,x)$.
    \item \textbf{Triangle Inequality}: $d(x,z) \le d(x,y) + d(y,z)$.
\end{enumerate}
The pair $(X,d)$ is called a \textbf{metric space}. The elements of $X$ are called \textbf{points}, and $d(x,y)$ is called the distance between $x$ and $y$.
\end{definition}

\begin{example}[Real Line $\R$]
Define $d(x,y) = |x-y|$ for all $x,y \in \R$. Then $(\R, d)$ is the usual metric space of real numbers.
\end{example}

\begin{example}[Euclidean Space $\R^n$]
For $x = (x_1, \dots, x_n)$ and $y = (y_1, \dots, y_n)$ in $\R^n$, define:
\[
d_2(x,y) = \left( \sum_{i=1}^n (x_i - y_i)^2 \right)^{1/2}.
\]
Then $(\R^n, d_2)$ is the $n$-dimensional Euclidean metric space.
\end{example}

\begin{example}[Taxicab and Maximum Metrics on $\R^n$]
On $\R^n$, we can also define:
\begin{align*}
\text{\textbf{Taxicab / $\ell_1$ Metric:}} \qquad d_1(x,y) &= \sum_{i=1}^n |x_i - y_i|, \\
\text{\textbf{Maximum / $\ell_\infty$ Metric:}} \qquad d_\infty(x,y) &= \max_{1 \le i \le n} |x_i - y_i|.
\end{align*}
Both $(\R^n, d_1)$ and $(\R^n, d_\infty)$ are metric spaces.
\end{example}

\begin{example}[Discrete Metric Space]
Let $X$ be any non-empty set. Define $d: X \times X \to \R$ by:
\[
d(x,y) = \begin{cases} 0 & \text{if } x = y, \\ 1 & \text{if } x \neq y. \end{cases}
\]
Then $d$ is called the \textbf{discrete metric}, and $(X,d)$ is a discrete metric space.
\end{example}

\begin{example}[Sequence Space $\ell_p$ ($1 \le p < \infty$)]
The set $\ell_p$ consists of all infinite sequences $x = (x_n)_{n=1}^\infty$ of real (or complex) numbers such that $\sum_{n=1}^\infty |x_n|^p < \infty$. The metric is defined by:
\[
d_p(x,y) = \left( \sum_{n=1}^\infty |x_n - y_n|^p \right)^{1/p}.
\]
By Minkowski's inequality, $d_p$ satisfies the triangle inequality.
\end{example}

\begin{example}[Space of Bounded Sequences $\ell_\infty$]
The set $\ell_\infty$ consists of all bounded sequences $x = (x_n)_{n=1}^\infty$ with metric:
\[
d_\infty(x,y) = \sup_{n \ge 1} |x_n - y_n|.
\]
\end{example}

\begin{example}[{Function Space $C[a,b]$}]
Let $C[a,b]$ be the set of all continuous real-valued functions defined on a closed bounded interval $[a,b]$. The uniform metric is defined by:
\[
d_\infty(f,g) = \max_{x \in [a,b]} |f(x) - g(x)|.
\]
We can also define the $L_1$ metric on $C[a,b]$ by:
\[
d_1(f,g) = \int_a^b |f(x) - g(x)|\,dx.
\]
\end{example}


\section{Open Spheres and Closed Spheres}
\label{sec:spheres}

\begin{definition}[Spheres and Balls]
Let $(X,d)$ be a metric space, $x_0 \in X$, and $r > 0$.
\begin{enumerate}[label=\rm{(\roman*)}]
    \item The \textbf{open sphere} (or open ball) centered at $x_0$ with radius $r$ is:
    \[
    B(x_0, r) = \{x \in X : d(x, x_0) < r\}.
    \]
    \item The \textbf{closed sphere} (or closed ball) centered at $x_0$ with radius $r$ is:
    \[
    \overline{B}(x_0, r) = \{x \in X : d(x, x_0) \le r\}.
    \]
    \item The \textbf{sphere} centered at $x_0$ with radius $r$ is:
    \[
    S(x_0, r) = \{x \in X : d(x, x_0) = r\}.
    \]
\end{enumerate}
\end{definition}

% TikZ Figure comparing Unit Balls in R^2
\begin{figure}[htbp]
\centering
\begin{subfigure}[b]{0.31\textwidth}
\centering
\begin{tikzpicture}[scale=1.5]
    % Axes
    \draw[->, thick, >=stealth] (-1.3,0) -- (1.3,0) node[right] {$x$};
    \draw[->, thick, >=stealth] (0,-1.3) -- (0,1.3) node[above] {$y$};
    % Ball d_1 (Rhombus)
    \filldraw[fill=mainblue!20, draw=mainblue, thick] (1,0) -- (0,1) -- (-1,0) -- (0,-1) -- cycle;
    \node[mainblue] at (0.4, 0.4) {$B_{d_1}(0,1)$};
    \draw[dotted] (1,0) node[below right] {$1$};
    \draw[dotted] (0,1) node[above left] {$1$};
\end{tikzpicture}
\caption{$d_1$ metric: $|x|+|y| < 1$}
\label{fig:ball_d1}
\end{subfigure}
\hfill
\begin{subfigure}[b]{0.31\textwidth}
\centering
\begin{tikzpicture}[scale=1.5]
    % Axes
    \draw[->, thick, >=stealth] (-1.3,0) -- (1.3,0) node[right] {$x$};
    \draw[->, thick, >=stealth] (0,-1.3) -- (0,1.3) node[above] {$y$};
    % Ball d_2 (Circle)
    \filldraw[fill=darkred!20, draw=darkred, thick] (0,0) circle (1);
    \node[darkred] at (0.4, 0.4) {$B_{d_2}(0,1)$};
    \draw[dotted] (1,0) node[below right] {$1$};
    \draw[dotted] (0,1) node[above left] {$1$};
\end{tikzpicture}
\caption{$d_2$ metric: $\sqrt{x^2+y^2} < 1$}
\label{fig:ball_d2}
\end{subfigure}
\hfill
\begin{subfigure}[b]{0.31\textwidth}
\centering
\begin{tikzpicture}[scale=1.5]
    % Axes
    \draw[->, thick, >=stealth] (-1.3,0) -- (1.3,0) node[right] {$x$};
    \draw[->, thick, >=stealth] (0,-1.3) -- (0,1.3) node[above] {$y$};
    % Ball d_infty (Square)
    \filldraw[fill=darkgreen!20, draw=darkgreen, thick] (-1,-1) rectangle (1,1);
    \node[darkgreen] at (0.4, 0.4) {$B_{d_\infty}(0,1)$};
    \draw[dotted] (1,0) node[below right] {$1$};
    \draw[dotted] (0,1) node[above left] {$1$};
\end{tikzpicture}
\caption{$d_\infty$ metric: $\max(|x|,|y|) < 1$}
\label{fig:ball_dinfty}
\end{subfigure}
\caption{Unit open balls centered at the origin $(0,0)$ in $\R^2$ under three distinct metrics.}
\label{fig:unit_balls_comparison}
\end{figure}

\begin{theorem}
In any metric space $(X,d)$, every open sphere $B(x_0, r)$ is an open set, and every closed sphere $\overline{B}(x_0, r)$ is a closed set.
\end{theorem}


\section{Neighbourhoods and Open Sets}
\label{sec:open_sets}

\begin{definition}[Neighbourhood]
Let $(X,d)$ be a metric space and $x \in X$. A subset $N \subseteq X$ is called a \textbf{neighbourhood} of $x$ if there exists $r > 0$ such that $B(x,r) \subseteq N$.
\end{definition}

\begin{definition}[Open Set]
A subset $G \subseteq X$ is called an \textbf{open set} if $G$ is a neighbourhood of each of its points, i.e., for every $x \in G$, there exists $r_x > 0$ such that $B(x, r_x) \subseteq G$.
\end{definition}

% TikZ Figure showing an Open Set and an interior open ball
\begin{figure}[htbp]
\centering
\begin{tikzpicture}[scale=1.2]
    % Smooth region representing open set G
    \draw[thick, fill=mainblue!10, draw=mainblue, smooth cycle] 
        plot coordinates {(0,0) (2,1) (4,0.5) (4.5,-1.5) (2.5,-2.5) (-0.5,-1.5)};
    \node[mainblue] at (3.8, 0.1) {$G$};
    
    % Point x and ball B(x,r)
    \coordinate (X) at (1.8, -0.8);
    \filldraw[black] (X) circle (1.5pt) node[above] {$x$};
    \draw[dashed, draw=darkred, fill=darkred!20] (X) circle (0.9);
    \draw[->, >=stealth, darkred] (X) -- +(0.9, 0) node[midway, above] {$r$};
    \node[darkred] at (2.2, -1.4) {$B(x,r) \subseteq G$};
\end{tikzpicture}
\caption{Geometric characterisation of an open set $G$: for any $x \in G$, an open sphere $B(x,r)$ is completely contained in $G$.}
\label{fig:open_set_illustration}
\end{figure}

\begin{theorem}[Topology of Open Sets]
In any metric space $(X,d)$:
\begin{enumerate}[label=\rm{(\roman*)}]
    \item The empty set $\emptyset$ and the whole space $X$ are open sets.
    \item The union of an arbitrary family of open sets is open.
    \item The intersection of a finite family of open sets is open.
\end{enumerate}
\end{theorem}


\section{Interior Points, Closed Sets, Limit Points, and Closure}
\label{sec:topological_notions}

\begin{definition}[Interior Point and Interior]
Let $A \subseteq X$. A point $x \in A$ is an \textbf{interior point} of $A$ if there exists $r > 0$ such that $B(x,r) \subseteq A$. The set of all interior points of $A$ is called the \textbf{interior} of $A$, denoted by $\Int(A)$ or $A^\circ$.
\end{definition}

\begin{theorem}
For any $A \subseteq X$, $\Int(A)$ is the largest open set contained in $A$. Furthermore, $A$ is open if and only if $\Int(A) = A$.
\end{theorem}

\begin{definition}[Closed Set]
A subset $F \subseteq X$ is called a \textbf{closed set} if its complement $F^c = X \setminus F$ is an open set.
\end{definition}

\begin{definition}[Limit Point and Derived Set]
A point $x \in X$ is called a \textbf{limit point} (or accumulation point) of $A \subseteq X$ if every open sphere $B(x,r)$ contains at least one point of $A$ other than $x$, i.e.,
\[
(B(x,r) \setminus \{x\}) \cap A \neq \emptyset \quad \forall r > 0.
\]
The set of all limit points of $A$ is called the \textbf{derived set} of $A$, denoted by $A'$.
\end{definition}

\begin{theorem}
A set $F \subseteq X$ is closed if and only if it contains all its limit points, i.e., $A' \subseteq F$.
\end{theorem}

\begin{definition}[Closure of a Set]
The \textbf{closure} of $A \subseteq X$, denoted by $\overline{A}$ or $\Cl(A)$, is defined as $\overline{A} = A \cup A'$.
\end{definition}

\begin{theorem}
The closure $\overline{A}$ is the smallest closed set containing $A$. Moreover, $A$ is closed if and only if $\overline{A} = A$.
\end{theorem}

\begin{definition}[Boundary Points]
A point $x \in X$ is a \textbf{boundary point} of $A \subseteq X$ if every open sphere $B(x,r)$ contains at least one point of $A$ and at least one point of $X \setminus A$. The set of all boundary points of $A$ is the \textbf{boundary} $\Bd(A) = \overline{A} \cap \overline{X \setminus A}$.
\end{definition}


\section{Distance Between Sets, Subspaces, and Product Spaces}
\label{sec:subspace_product}

\begin{definition}[Distance and Diameter]
Let $(X,d)$ be a metric space, $A, B \subseteq X$, and $x \in X$.
\begin{enumerate}[label=\rm{(\roman*)}]
    \item Distance from $x$ to $A$: $d(x,A) = \inf \{d(x,a) : a \in A\}$.
    \item Distance between $A$ and $B$: $d(A,B) = \inf \{d(a,b) : a \in A, b \in B\}$.
    \item Diameter of $A$: $\diam(A) = \sup \{d(a_1, a_2) : a_1, a_2 \in A\}$.
\end{enumerate}
\end{definition}

\begin{definition}[Subspace Metric]
Let $(X,d)$ be a metric space and $Y \subseteq X$ a non-empty subset. The restriction $d_Y = d|_{Y \times Y}$ is a metric on $Y$, and $(Y, d_Y)$ is called a \textbf{metric subspace} of $(X,d)$.
\end{definition}
